% =============================================================================
% Project : THE_PHILOSOPHY_OF_VIABILITY
% Book    : Book I — Ontology of Viability
% Book ID : VS.B01
%
% File    : 4_ASSETS/tables/VS.B01.P01.C02.S03.SS10.P001.TL01.conditions_of_an_invariance_claim.tex
%
% Section / Chapter / Asset:
%   Chapter 02 — Invariance
%   Section 03 — Defining Invariants
%   Table TL01 — Conditions of an Invariance Claim
%
% Asset ID:
%   VS.B01.P01.C02.S03.SS10.P001.TL01.conditions_of_an_invariance_claim
%
% Purpose :
%   To distinguish the principal analytical elements and validity
%   conditions required for formulating and evaluating an invariance
%   claim.
%
% Conceptual Policy:
%   1. The table presents analytical requirements, not ontological
%      components.
%   2. Specification, scope, and ontological grounding have distinct
%      logical functions.
%   3. The table does not presuppose that an invariant exists.
%   4. The requirements are not presented as a universal formal
%      algorithm.
%
% =============================================================================

\begin{table}[htbp]
    \centering

    \caption{
        Analytical Conditions for Formulating and Evaluating
        an Invariance Claim
    }
    \label{tab:conditions-of-an-invariance-claim}

    \vspace{0.5em}

    \small

    \renewcommand{\arraystretch}{1.45}

    \begin{tabular}{
        |p{0.23\textwidth}
        |p{0.65\textwidth}|
    }
        \hline

        \textbf{Analytical Element}
        &
        \textbf{Question and Function}
        \\ \hline

        \textbf{Subject}
        &
        What determination, relation, structure, or other
        identifiable feature is claimed to be preserved?
        \\ \hline

        \textbf{Transformation}
        &
        Under which specified change or class of changes
        is preservation being examined?
        \\ \hline

        \textbf{Domain}
        &
        Within which range of configurations, conditions,
        or cases does the transformation and the claim apply?
        \\ \hline

        \textbf{Respect}
        &
        In what particular regard is the subject claimed
        to remain preserved?
        \\ \hline

        \textbf{Preservation Criterion}
        &
        What standard determines whether the specified
        subject remains unchanged in the relevant respect?
        \\ \hline

        \textbf{Transformation Admissibility}
        &
        Which transformations legitimately belong to the
        specified class against which preservation is assessed?
        \\ \hline

        \textbf{Scope of Validity}
        &
        What are the limits within which the preservation
        claim is justified, and beyond which it cannot
        be generalized without further argument?
        \\ \hline

        \textbf{Ontological Grounding}
        &
        What justifies interpreting the established
        preservation as a feature of the relevant existent
        or relation rather than merely as a formal result?
        \\ \hline

    \end{tabular}

    \normalsize

    \vspace{0.4em}

    \begin{minipage}{0.92\textwidth}
        \footnotesize
        \textit{Note.}
        The elements listed above have different analytical
        functions. They specify the content, applicability,
        limits, and ontological interpretation of an
        invariance claim. They are not independent
        ontological components and do not imply that
        an invariant necessarily exists.
    \end{minipage}

\end{table}

% =============================================================================
% END OF TABLE TL01
% =============================================================================